\chapter{Propriospinal Myoclonus at Sleep Onset}
\index{Propriospinal Myoclonus at Sleep Onset}
\label{sec:Propriospinal Myoclonus at Sleep Onset}

\section{Overview}
Propriospinal myoclonus at sleep onset is a rare sleep-related movement disorder characterized by sudden, brief, involuntary jerks of the axial muscles occurring during the transition from wakefulness to sleep. This condition is among the most common mental health disorders worldwide. It affects people across all age groups, socioeconomic backgrounds, and geographic regions.
\section{Causes}
This condition happens because of myoclonus originating in the spinal cord and propagating rostrally and caudally via propriospinal pathways.     
\section{Risk Factors}

\begin{itemize}[noitemsep]
  \item Genetic predisposition and family history
  \item Advancing age
  \item Lifestyle factors (diet, exercise, smoking, alcohol)
  \item Environmental exposures
  \item Underlying medical conditions
  \item Medication use
\end{itemize}
\section{Signs and Symptoms}

\subsection{Common symptoms}
\begin{itemize}[noitemsep]
  \item You may experience jerks triggered by voluntary attempts to relax or fall asleep, severely delaying sleep onset and resulting in sleep-onset insomnia. 
  \end{itemize}

\subsection{Emergency warning signs}
\begin{itemize}[noitemsep]
  \item Severe or worsening symptoms of propriospinal myoclonus at sleep onset require immediate medical evaluation
\end{itemize}
\section{Progression and Stages}
The condition may progress through different stages of severity over time. Early stages often involve mild or subtle symptoms that may be overlooked. Moderate stages involve more noticeable symptoms affecting daily function. Advanced stages may involve significant impairment and complications.
\section{Complications}
\begin{itemize}[noitemsep]
  \item Disease progression leading to worsening symptoms
  \item Secondary infections or injuries
  \item Side effects of treatment
  \item Reduced quality of life
  \item Emotional and psychological impact
\end{itemize}
\section{Diagnosis}
To make the diagnosis, doctors look for polygraphic EMG recording across multiple spinal segments and polysomnography, with differentiation from hypnic jerks (physiologic sleep starts) and epileptic myoclonus. Comprehensive medical history and physical examination Laboratory tests to assess organ function and rule out other conditions Imaging studies to visualize affected structures Specialized testing depending on the affected system Consultation with relevant specialists Monitoring of disease progression over time

\begin{itemize}[noitemsep]
  \item Comprehensive medical history and physical examination
  \item Laboratory tests to assess organ function
  \item Imaging studies to visualize affected structures
  \item Specialized testing depending on the affected system
  \item Consultation with relevant specialists
\end{itemize}
\section{Prevention}
\begin{itemize}[noitemsep]
  \item Maintain a healthy lifestyle with balanced diet and exercise
  \item Avoid known risk factors
  \item Attend regular medical check-ups for early detection
  \item Follow recommended screening guidelines
\end{itemize}
\section{Treatment Options}
\begin{itemize}[noitemsep]
  \item Treatment typically involves challenging
  \item Clonazepam, levetiracetam, and other anticonvulsants have been used with variable success.
  \item Medications to manage symptoms and address underlying mechanisms
  \item Lifestyle modifications and self-care measures
  \item Physical therapy and rehabilitation
  \item Surgical intervention when indicated
  \item Regular monitoring and follow-up care
\end{itemize}
\section{Living with It}
\begin{itemize}[noitemsep]
  \item Follow your treatment plan as prescribed
  \item Maintain regular follow-up with your healthcare provider
  \item Make necessary lifestyle adjustments
  \item Seek support from family, friends, or support groups
  \item Stay informed about your condition
\end{itemize}
\section{Outlook}
\begin{itemize}[noitemsep]
  \item The outlook varies from person to person
  \item Some patients experience spontaneous improvement, while others require long-term pharmacotherapy.
\end{itemize}
